\centerline{\bf \Huge Производные}

\problem{\textbf{1}} Докажите с помощью определения производной:\\
а) $f(x)=x^2$ => $f'(x)=2x$\\
б) $f(x)=x^3$ => $f'(x)=3x^2$\\
в) $f(x)=\dfrac{1}{x}$ => $f'(x)=-\dfrac{1}{x^2}$\\
г) $f(x)=\sqrt{x}$ => $f'(x)=\dfrac{1}{2\sqrt{x}}$\\
д) $f(x)=x^n$ => $f'(x)=nx^{n-1}$, если $n \in \boldsymbol{N}$

\problem{\textbf{2}} Найдите первую производную:\\
а) $y=6+x+3 x^2-\sin x-2 \sqrt[3]{x}+\dfrac{1}{x^2}-11 \operatorname{ctg} x$\\
б) $y=5 \ln x+\dfrac{2}{\sqrt[5]{x^7}}+\operatorname{arctg} x-\lg 20 \cdot 2^x+\operatorname{tg} 3$\\
в) $y=\left(x^2+7 x-1\right) \log _3 x$\\
г) $y=\dfrac{3^x+5}{\cos x}$\\
д) $y=(2 x+1)^5$\\
е) $y=\dfrac{1}{\left(x^2-1\right)^7}$\\
ж) $y=\sqrt{\operatorname{arcctg} x}$\\
з) $y=\sqrt[3]{x^2+\operatorname{tg} x+15}$\\
и) $y=\dfrac{5}{\sqrt[5]{x+\ln x}}$\\
к) $y=-\dfrac{1}{\cos x}$\\
л) $y=7^{\sin ^2 x}$\\
м) $y=\ln ^2(2 x-1)$\\
н) $y=-2 x e^{3 x}+\left(x^2-4 x+3\right) \sin 7 x$\\
о) $y=\sqrt{x^2+4} \cdot \ln (\sin x)$

\problem{\textbf{3}} а) Найдите точку минимума и наименьшее значение функции $y=\sqrt{x^2-28 x+211}-49$\\
б) Найдите точку максимума и наибольшее значение функции $y=5\log _5\left(-19+12 x-x^2\right)$\\
в) Найдите точку минимума и наименьшее значение функции $y=24\cdot3^{x^2+14 x+71}+1$.\\
г) Найдите точку максимума и наибольшее значение функции $y=(x+9)^2(x-2)-3$\\
д) Найдите точку минимума и наименьшее значение функции $y=x^3+12 x^2+36 x+51$\\
е) Найдите точку максимума и наибольшее значение функции $y=-\dfrac{x}{x^2+144}$\\
ж) Найдите точку минимума и наименьшее значение функции $y=\left(x^2-5 x-5\right) e^{7-x}$

\newpage


\problem{\textbf{4}} Найдите наименьшее и наибольшее значение функции на отрезке:\\
а) $y=7+12 x-x^3$,$[-3 ; 4]$\\
б) $y=\dfrac{x^2+25}{x}$,$[-10 ; -1]$\\
в) $(x-2)^2 e^{x-6}$,$[-1 ; 5]$\\
г) $y=9 x-\ln (9 x)+3$,$\left[\dfrac{1}{18} ; \dfrac{5}{18}\right]$

\problem{\textbf{5}} а) $y=x^3-7x^{20}+19x-2$, $y^{(21)}=?$\\
б) $y=-101x^{100}+5x^{99}-2x^{49}$, $y^{(99)}=?$ \\
в) $y=\ln (2 x-1)$, $y^{(2026)}=?$\\
г) $y=\sqrt{x-76}$, $y^{(67)}=?$

\problem{\textbf{7}} Два велосипедиста равномерно движутся по взаимно перпендикулярным дорогам по направлению к перекрестку этих дорог. Один из них движется со скоростью 40 км/ч и находится на расстоянии 5 км от перекрестка, второй движется со скоростью 30 км/ч и находится на расстоянии 3 км от перекрестка. Через сколько минут расстояние между велосипедистами станет наименьшим? Каково будет это наименьшее расстояние? Считайте, что перекресток не T-образный, обе дороги продолжаются за перекрестком.

\problem{\textbf{8}} Антон является владельцем двух заводов в разных городах. На заводах производится абсолютно одинаковые товары при использовании одинаковых технологий. Если рабочие на одном из заводов трудятся суммарно $t^2$ часов в неделю, то за эту неделю они производят $t$ единиц товара.За каждый час работы на заводе, расположенном в первом городе, Антон платит рабочему 250 рублей, а на заводе, расположенном во втором городе, - 200 рублей. Антон готов выделять 900000 рублей в неделю на оплату труда рабочих. Какое наибольшее количество единиц товара можно произвести за неделю?

\problem{\textbf{9}} Зависимость количества $Q$ (в шт., $0 \leqslant Q \leqslant 20000$ ) купленного у фирмы товара от цены $P$ (в руб. за шт.) выражается формулой $Q=20000-P$. Затраты на производство $Q$ единиц товара составляют $6000 Q+4000000$ рублей. Кроме затрат на производство, фирма должна платить налог $t$ рублей $(0<t<10000)$ с каждой произведённой единицы товара. Таким образом, прибыль фирмы составляет $P Q-6000 Q-4000000-t Q$ рублей, а общая сумма налогов, собранных государством, равна $t Q$ рублей. Фирма производит такое количество товара, при котором её прибыль максимальна. При каком значении $t$ общая сумма налогов, собранных государством, будет максимальной?\\


\centerline{\textbf{180 минут}}